\documentclass[titlepage,a4paper]{jsarticle}
\usepackage{../../../sty/import}% 各種パッケージインポート
\usepackage{../../../sty/title}% タイトルページの変更

%% タイトルページの変数
% レポートタイトル
\title{第一回授業課題}
% 提出日
\expdate{\today}
% 科目名
\subject{原子力材料と核燃料}
% 分野
\class{情報経営システム工学分野}
% 学年
\grade{M1}
% 学籍番号
\mynumber{24336488}
% 記述者
\author{本間三暉}

\begin{document}
% titleページ作成
\maketitle

\section{原子炉で使われている材料を一つ挙げ，その特徴を述べよ．}
\subsection{原子力材料}
ジルカロイ
\subsection{特徴}
ジルコニウム元素を素材にした，錫，鉄，ニッケル，クローム等の小量を含む合金である．
ジルカロイは軽量で高強度，耐食性に優れているため，原子炉の構造材料として使用され，主に燃料被覆管や構造部材に利用されている．
また，ジルカロイは中性子吸収が少ないため，核反応を妨げずに効率的なエネルギー生成が可能である．

% 参考文献
\begin{thebibliography}{99}
\bibitem{1}原子力百科事典 ATOMICA \url{https://atomica.jaea.go.jp/dic/detail/dic_detail_394.html}
\bibitem{2}一般社団法人 日本機械学会 \url{https://www.jsme.or.jp/jsme-medwiki/doku.php?id=12:1006006}

\end{thebibliography}

\end{document}